\documentclass{scalatekids-article}
\begin{document}
\lfoot{Specifica Tecnica 1.0.0}
\newgeometry{top=3.5cm}
\begin{titlepage}
  \begin{center}
    \begin{center}
      \includegraphics[width=10cm]{sklogo.png}
    \end{center}
    \vspace{1cm}
    \begin{Huge}
      \begin{center}
        \textbf{Specifica Tecnica}
      \end{center}
    \end{Huge}
    \vspace{11pt}
    \bgroup
    \def\arraystretch{1.3}
    \begin{tabular}{r|l}
      \multicolumn{2}{c}{\textbf{Informazioni sul documento}} \\
      \hline
      \setbox0=\hbox{0.0.1\unskip}\ifdim\wd0=0pt
      \\
      \else
      \textbf{Versione} & 1.0.0\\
      \fi
      \textbf{Redazione} & \multiLineCell[t]{}\\
      \textbf{Verifica} & \multiLineCell[t]{}\\
      \textbf{Approvazione} & \multiLineCell[t]{}\\
      \textbf{Uso} & Esterno\\
      \textbf{Lista di Distribuzione} & \multiLineCell[t]{ScalateKids\\Prof. Tullio Vardanega\\Prof. Riccardo Cardin}\\
    \end{tabular}
    \egroup
    \vspace{22pt}
  \end{center}
\end{titlepage}
\restoregeometry
\clearpage
\pagenumbering{Roman}
\setcounter{page}{1}
\begin{flushleft}
  \vspace{0cm}
  {\large\bfseries Diario delle modifiche \par}
\end{flushleft}
\vspace{0cm}
\begin{center}
  \begin{tabular}{| l | l | l | l | p{5cm} |}
    \hline
    Versione & Autore & Ruolo & Data & Descrizione \\
    \hline
    0.2.1 & Marco Boseggia & Progettista & 2016-03-21 & Stesura sezioni 4.6 4.7 4.8 4.9\\
    \hline
    0.2.0 & Michael Munaro & Verificatore & 2016-03-21 & Verifica sottosezioni 3.1 3.2 4.1 4.2 4.3 4.4 4.5\\
    \hline
    0.1.2 & Andrea Giacomo Baldan & Progettista & 2016-03-17 & Stesura sottosezioni 4.1 4.2 4.3 4.4 4.5\\
    \hline
    0.1.1 & Andrea Giacomo Baldan & Progettista & 2016-03-17 & Stesura sottosezioni 3.1 3.2\\
    \hline
    0.1.0 & Michael Munaro & Verificatore & 2016-03-16 & Verifica sezioni 1 e 2\\
    \hline
    0.0.2 & Marco Boseggia & Progettista & 2016-02-26 & Stesura sezioni 1 e 2\\
    \hline
    0.0.1 & Andrea Giacomo Baldan & Amministratore & 2016-02-26 & Creazione scheletro del documento\\
    \hline
  \end{tabular}
\end{center}
\tableofcontents
\newpage
\pagenumbering{arabic}
\section{Sommario}
\subsection{Scopo del documento}
Il seguente documento ha lo scopo di descrivere la progettazione ad alto livello che il gruppo \textit{ScalateKids} ha scelto per il progetto \textbf{ActorBase}.\\
Sarà descritta l'\gloss{architettura} generale del progetto, in particolare la scelta dei \gloss{design pattern} e delle componenti che andranno a comporre il software.
\prodPurpose
\glossExpl
\subsection{Riferimenti}
\subsubsection{Normativi}
\begin{itemize}
\item\textbf{Capitolato d'appalto C1:} \textit{Actorbase: a NoSQL DB based on the Actor model;}\\
  \url{http://www.math.unipd.it/~tullio/IS-1/2015/Progetto/C1.pdf}
\item\textbf{Norme di Progetto:} \href{run:../Interni/NormeDiProgetto\_v1.0.0.pdf}{Norme di Progetto v1.0.0.}
\end{itemize}
\subsubsection{Informativi}
\begin{itemize}
\item\textbf{Piano di Progetto:} \href{run:./PianoDiProgetto\_v1.0.0.pdf}{Piano di Progetto v1.0.0}
\item\textbf{Dispense fornite dall'insegnamento Ingegneria del Software mod. A:}\\
  \url{http://www.math.unipd.it/~tullio/IS-1/2015/}
\item\textbf{Documentazione di \gloss{Akka} su \gloss{Scala}:}
  \url{http://doc.akka.io/docs/akka/2.4.2/scala.html}
\end{itemize}
\newpage
\section{Tecnologie utilizzate}
In questa sezione saranno descritte le tecnologie scelte per lo sviluppo del progetto \textbf{ActorBase} e le motivazioni che ci hanno spinto a sceglierle.
\subsection{Akka}
La scelta della libreria \gloss{Akka} è stata dettata dal capitolato. Tuttavia l'avremmo scelta comunque per i seguenti motivi:
\begin{itemize}
\item\textbf{Implementazione del \gloss{modello ad attori}:} \gloss{Akka} rende semplice la creazione e l'utilizzo degli \gloss{attori}. Grazie a ciò viene naturale la distribuibilità del progetto e l'asincronia tra i processi;
\item\textbf{Tolleranza agli errori:} \gloss{Akka} implementa un sistema di gerarchia di supervisori. Questa gerarchia consente ai supervisori di un \gloss{attore}, nel caso in cui quest'ultimo lanci un'eccezione, di mandarlo in \gloss{crash} e farlo poi ripartire. In questo modo il sistema è resiliente agli errori;
\item\textbf{Persistenza:} Ogni messaggio ricevuto da ciascun \gloss{attore} può essere salvato in maniera tale che al riavvio esso possa riprendere da dov'era prima del suo spegnimento.
\end{itemize}
\subsection{Scala}
Il capitolato richiede l'utilizzo di un linguaggio di programmazione a scelta tra \gloss{Scala} e \gloss{Java}. Abbiamo scelto \gloss{Scala} per i seguenti motivi:
\begin{itemize}
\item\textbf{\gloss{Programmazione funzionale}:} Questa tecnica di programmazione evita che ci siano degli effetti collaterali tra funzioni, rendendole quindi \gloss{thread-safe};
\item\textbf{Implementazione di \gloss{Akka}:} \gloss{Akka} risulta più facilmente implementabile in \gloss{Scala} che in \gloss{Java}. Questo perchè la \gloss{programmazione funzionale} con i suoi paradigmi si integra meglio con il \gloss{modello ad attori};%%TODO completare lista
\end{itemize}

\section{Descrizione architettura}

\subsection{Metodo e formalismo di specifica}

L'approccio alla progettazione è stato meet-in-the-middle, la prima suddivisione ad alto livello è stata tra sistema come struttura ad attori e
componenti esterne per utilizzare la base di dati (operazioni CRUD)

\subsection{Architettura generale}

Macroscopicamente il sistema si suddivide in due componenti principali, assimilabili ad un paradigma Client-Server, il sistema ad attori
dove risiedono la struttura della base di dati ed un server TCP pronto a ricevere comandi dall'esterno rappresentano la componente server,
la Command Line Interface (CLI) rappresenta la componente client mediante la quale è possibile inviare comandi al sistema e ricevere l'output
prodotto da essi, essa farà uso della componente \gloss{driver} per comunicare con il sistema lato server mediante un protocollo di comunicazione
appositamente creato.

\begin{figure}[H]
  \begin{center}
    \includegraphics[width=0.2\textwidth,keepaspectratio]{RP/RP-ClientServer.png}
    \caption*{Diagramma architettura generale}
  \end{center}
\end{figure}

% TODO:
% \begin{figure}
%   \centering

%   \caption{Architettura generale - vista package}
%   \label{fig:architetturaGeneralePackage}
% \end{figure}

L'intero sistema è rappresentabile mediante una variante del design pattern
\gloss{MVC}, le tre componenti principali si prestano naturalmente
all'assegnazione delle tre parti che formano il design pattern, con la
differenza rispetto ad un approccio più classico, che l'aggiornamento della
parte \gloss{view} avviene sempre e comunque passando per la parte controller,
ciò è dovuto alla natura client-server del prodotto, che richiede un protocollo
di comunicazione ben preciso tra le due parti. Il sistema formato dai package
dedicati agli attori \gloss{akka} e dal server \gloss{TCP} per la comunicazione
con l'esterno sono assimilabili alla parte \gloss{model}, rappresenta la
\gloss{business logic} di \textbf{Actorbase}. La \gloss{CLI} è assimilabile alla
parte \gloss{view}, unico punto di interazione con l'utente, si occupa di
processare i comandi in input e restituire l'output opportunamente formattato,
la componente \gloss{driver} fa da ponte di comunicazione tra \gloss{CLI} e
server, occupandosi della ``traduzione'' dei comandi input secondo il protocollo
di comunicazione scelto, rappresenta la parte \gloss{controller}.

\section{Componenti}

\subsection{actorbase}

% TODO:
% \begin{figure}
%   \centering

%   \caption{Componente Actorbase}
%   \label{fig:componenteActorbase}
% \end{figure}

\subsubsection{Descrizione}

La componente principale rappresentata dal package globale \textbf{actorbase},
contiene le parti tipiche del design pattern \gloss{MVC}.

\subsubsection{Package contenuti}

\begin{itemize}
\item actorbase::model;
\item actorbase::controller;
\item actorbase::view.
\end{itemize}

Rappresentano le tre macro-componenti del sistema.

\subsection{actorbase::model}

\subsubsection{Descrizione}

\gloss{Package} per la componente \gloss{Model} dell'architettura \gloss{MVC}, rappresenta anche la componente system del prodotto.

\subsubsection{Package contenuti}

\begin{itemize}
\item system.
\end{itemize}

\subsection{actorbase::model::system}

\subsubsection{Descrizione}

\gloss{Package} per la componente system del prodotto.

\subsubsection{Package contenuti}

\begin{itemize}
\item server;
\item actors.
\end{itemize}

\subsection{actorbase::controller}

\subsubsection{Descrizione}

\gloss{Package} per la componente \gloss{Controller} dell'architettura \gloss{MVC}.

\subsubsection{Package contenuti}

\begin{itemize}
\item driver.
\end{itemize}

\subsection{actorbase::controller::driver}

\subsubsection{Descrizione}

\gloss{Package} per la componente \gloss{driver} del prodotto.

\subsubsection{Package contenuti}

\begin{itemize}
\item client;
\item controls;
\item data.
\end{itemize}

\subsection{actorbase::view}

\subsubsection{Descrizione}

\gloss{Package} per la componente \gloss{View} dell'architettura \gloss{MVC}.

\subsubsection{Package contenuti}

\begin{itemize}
\item cli.
\end{itemize}

\subsection{actorbase::view::cli}

\subsubsection{Descrizione}

\gloss{Package} per la componente \gloss{cli} del prodotto.

\subsubsection{Package contenuti}

\begin{itemize}
\item actorbaseCli;
\item commands.
\end{itemize}

\subsection{actorbase::view::cli::actorbaseCli}

\subsubsection{Descrizione}

Componente parte della \gloss{View} per la gestione dei comandi e dell'output prodotto verso l'utente.

\subsubsection{Interazioni con altre componenti}

\begin{itemize}
\item actorbase::view::cli::commands;
\item SpringShell.
\end{itemize}

\subsubsection{Classi}

\paragraph{actorbase::view::cli::actorbaseCli::PromptProvider}

\subparagraph{Descrizione}

Classe per la creazione del \gloss{prompt} dei comandi.

\subparagraph{Utilizzo}

Viene utilizzata per la creazione di un \gloss{prompt} per l'inserimento dei comandi e la visualizzazione dell'output.

\subparagraph{Classi ereditate}

\begin{itemize}
\item DefaultPromptProvider;
\item SpringShell.
\end{itemize}

\paragraph{actorbase::view::cli::actorbaseCli::HistoryFileNameProvider}

\subparagraph{Descrizione}

Classe per la gestione dello storico dei comandi \gloss{CLI}.

\subparagraph{Utilizzo}

Viene utilizzata per la creazione e gestione del sistema di storico dei comandi \gloss{CLI}.

\subparagraph{Classi ereditate}

\begin{itemize}
\item DefaultFilenameProvider;
\item SpringShell.
\end{itemize}

\paragraph{actorbase::view::cli::actorbaseCli::BannerProvider}

\subparagraph{Descrizione}

Classe per la creazione del \gloss{banner} d'ingresso alla \gloss{CLI}.

\subparagraph{Utilizzo}

Viene utilizzata per la creazione di un \gloss{banner} di presentazione all'ingresso della \gloss{CLI}.

\subparagraph{Classi ereditate}

\begin{itemize}
\item DefaultBannerProvider;
\item SpringShell.
\end{itemize}

\paragraph{actorbase::view::cli::actorbaseCli::CLICommands}

\subparagraph{Descrizione}

Classe per la gestione dei comandi della \gloss{CLI}.

\subparagraph{Utilizzo}

Viene utilizzata per la gestione mediante metodi appositi dei comandi della \gloss{CLI}, è responsabile della creazione e gestione dell'help sia generale che specifico per comando.

\subparagraph{Classi ereditate}

\begin{itemize}
\item CommandMarker.
\end{itemize}

\subparagraph{Relazioni con altre classi}

\begin{itemize}
\item actorbase::view::cli::commands::Command.
\end{itemize}

\paragraph{actorbase::view::cli::actorbaseCli::ResultViewer}

\subparagraph{Descrizione}

Classe per la gestione e la formattazione in output dei risultati ottenuti dall'esecuzione del comando ricevuto in input.

\subparagraph{Utilizzo}

Viene utilizzata per la gestione e la formattazione in output dei risultati ottenuti dal server mediante l'inserimento dei comandi.

\subsection{actorbase::view::cli::commands}

\subsubsection{Descrizione}

Componente parte della \gloss{View} per l'esecuzione del comando e l'interazione con il \gloss{Controller}.

\subsubsection{Interazioni con altre componenti}

\begin{itemize}
\item actorbase::view::cli::actorbaseCli;
\item actorbase::controller::driver::client.
\end{itemize}

\subsubsection{Classi}

\paragraph{actorbase::view::cli::commands::Command}

\subparagraph{Descrizione}

Interfaccia generica che rappresenta un comando.

\subparagraph{Utilizzo}

Viene utilizzata per offrire un'interfaccia comune a tutti i comandi.

\subparagraph{Relazioni con altre classi}

\begin{itemize}
\item actorbase::view::cli::actorbaseCli::CLICommands;
\item actorbase::view::cli::actorbaseCli::ResultViewer;
\item actorbase::controller::driver::client::ActorbaseClient.
\end{itemize}

\paragraph{actorbase::view::cli::commands::FindCommand}

\subparagraph{Descrizione}

Classe per il comando di ricerca.

\subparagraph{Utilizzo}

Viene utilizzata per mandare al \gloss{controller} il comando di ricerca con i parametri immessi dall'utente.

\subparagraph{Classi ereditate}

\begin{itemize}
\item actorbase::view::cli::commands::Command.
\end{itemize}

\paragraph{actorbase::view::cli::commands::LogoutCommand}

\subparagraph{Descrizione}

Classe per il comando di logout.

\subparagraph{Utilizzo}

Viene utilizzata per mandare al \gloss{controller} il comando di logout.

\subparagraph{Classi ereditate}

\begin{itemize}
\item actorbase::view::cli::commands::Command.
\end{itemize}

\paragraph{actorbase::view::cli::commands::ResetPasswordCommand}

\subparagraph{Descrizione}

Classe per il comando di reset della password.

\subparagraph{Utilizzo}

Viene utilizzata per mandare al \gloss{controller} il comando di reset della password dell'utente specificato in input.

\subparagraph{Classi ereditate}

\begin{itemize}
\item actorbase::view::cli::commands::Command.
\end{itemize}

\paragraph{actorbase::view::cli::commands::ListCommand}

\subparagraph{Descrizione}

Classe per il comando di visualizzazione dell'elenco dei nomi di tutte le collezioni presenti nel \gloss{database}.

\subparagraph{Utilizzo}

Viene utilizzata per mandare al \gloss{controller} il comando di visualizzazione dell'elenco dei nomi di tutte le collezioni presenti nel \gloss{database}.

\subparagraph{Classi ereditate}

\begin{itemize}
\item actorbase::view::cli::commands::Command.
\end{itemize}

\paragraph{actorbase::view::cli::commands::AddContributorCommand}

\subparagraph{Descrizione}

Classe per il comando di aggiunta collaboratore a una collezione.

\subparagraph{Utilizzo}

Viene utilizzata per mandare al \gloss{controller} il comando di aggiunta di un collaboratore a una collezione con il nome della \gloss{collezione} e lo username specificati in input.

\subparagraph{Classi ereditate}

\begin{itemize}
\item actorbase::view::cli::commands::Command.
\end{itemize}

\paragraph{actorbase::view::cli::commands::LoginCommand}

\subparagraph{Descrizione}

Classe per il comando di login.

\subparagraph{Utilizzo}

Viene utilizzata per mandare al \gloss{controller} il comando di login con le credenziali immesse dall'utente.

\subparagraph{Classi ereditate}

\begin{itemize}
\item actorbase::view::cli::commands::Command.
\end{itemize}

\paragraph{actorbase::view::cli::commands::DeleteCommand}

\subparagraph{Descrizione}

Classe per il comando di rimozione una \gloss{collezione}.

\subparagraph{Utilizzo}

Viene utilizzata per mandare al \gloss{controller} il comando di rimozione di una collezione con i parametri specificati in input.

\subparagraph{Classi ereditate}

\begin{itemize}
\item actorbase::view::cli::commands::Command.
\end{itemize}

\paragraph{actorbase::view::cli::commands::RemoveContributorCommand}

\subparagraph{Descrizione}

Classe per il comando di rimozione di un collaboratore da una collezione.

\subparagraph{Utilizzo}

Viene utilizzata per mandare al \gloss{controller} il comando di rimozione di un collaboratore da una collezione con con il nome della \gloss{collezione} e lo username specificati in input.

\subparagraph{Classi ereditate}

\begin{itemize}
\item actorbase::view::cli::commands::Command.
\end{itemize}

\paragraph{actorbase::view::cli::commands::RenameCommand}

\subparagraph{Descrizione}

Classe per il comando rinominazione di un item.

\subparagraph{Utilizzo}

Viene utilizzata per mandare al \gloss{controller} il comando di rinominazione di un item con i parametri specificati in input.

\subparagraph{Classi ereditate}

\begin{itemize}
\item actorbase::view::cli::commands::Command.
\end{itemize}

\paragraph{actorbase::view::cli::commands::AddUserCommand}

\subparagraph{Descrizione}

Classe per il comando di aggiunta di un utente.

\subparagraph{Utilizzo}

Viene utilizzata per mandare al \gloss{controller} il comando di aggiunta di un utente con i parametri specificati in input.

\subparagraph{Classi ereditate}

\begin{itemize}
\item actorbase::view::cli::commands::Command.
\end{itemize}

\paragraph{actorbase::view::cli::commands::RemoveCommand}

\subparagraph{Descrizione}

Classe per il comando di rimozione di un \gloss{item}.

\subparagraph{Utilizzo}

Viene utilizzata per mandare al \gloss{controller} il comando di rimozione di un \gloss{item} con i parametri immessi dall'utente.

\subparagraph{Classi ereditate}

\begin{itemize}
\item actorbase::view::cli::commands::Command.
\end{itemize}

\paragraph{actorbase::view::cli::commands::RemoveUserCommand}

\subparagraph{Descrizione}

Classe per il comando di rimozione di un utente.

\subparagraph{Utilizzo}

Viene utilizzata per mandare al \gloss{controller} il comando di rimozione di un utente con i parametri immessi dall'utente.

\subparagraph{Classi ereditate}

\begin{itemize}
\item actorbase::view::cli::commands::Command.
\end{itemize}

\paragraph{actorbase::view::cli::commands::ImportCommand}

\subparagraph{Descrizione}

Classe per il comando di importazione da un file \gloss{JSON}.

\subparagraph{Utilizzo}

Viene utilizzata per mandare al \gloss{controller} il comando di importazione da un file \gloss{JSON} con i parametri immessi dall'utente.

\subparagraph{Classi ereditate}

\begin{itemize}
\item actorbase::view::cli::commands::Command.
\end{itemize}

\paragraph{actorbase::view::cli::commands::InsertCommand}

\subparagraph{Descrizione}

Classe per il comando di inserimento di un \gloss{item} in una \gloss{collezione}.

\subparagraph{Utilizzo}

Viene utilizzata per mandare al \gloss{controller} il comando di inserimento di un \gloss{item} in una \gloss{collezione} con i parametri immessi dall'utente.

\subparagraph{Classi ereditate}

\begin{itemize}
\item actorbase::view::cli::commands::Command.
\end{itemize}

\paragraph{actorbase::view::cli::commands::CreateCommand}

\subparagraph{Descrizione}

Classe per il comando di creazione di una nuova \gloss{collezione}.

\subparagraph{Utilizzo}

Viene utilizzata per mandare al \gloss{controller} il comando di creazione di una nuova \gloss{collezione} con i parametri immessi dall'utente.

\subparagraph{Classi ereditate}

\begin{itemize}
\item actorbase::view::cli::commands::Command.
\end{itemize}

\paragraph{actorbase::view::cli::commands::ExportCommand}

\subparagraph{Descrizione}

Classe per il comando di esportazione in file \gloss{JSON}.

\subparagraph{Utilizzo}

Viene utilizzata per mandare al \gloss{controller} il comando di esportazione in file \gloss{JSON} con i parametri immessi dall'utente.

\subparagraph{Classi ereditate}

\begin{itemize}
\item actorbase::view::cli::commands::Command.
\end{itemize}

\section{Diagrammi attività}

\section{Design Pattern}

\subsection{Stime di fattibilità e di bisogno di risorse}

\section{Tracciamento}

\end{document}